\begin{center}
    \section*{\fontsize{20}{20}\selectfont Chapter 6}
\end{center}
\vspace{10mm}

\section{Conclusion}

This thesis demonstrates that graph neural networks constitute a viable and computationally efficient approach to brain tumour segmentation, challenging the prevailing assumption that volumetric 3D convolutions are necessary for competitive medical image analysis. The proposed framework achieves 92.92\% Dice coefficient through 5-fold ensemble aggregation on the BraTS 2021 dataset while delivering 6.9$\times$ faster inference and 156$\times$ fewer parameters compared to a standard 3D U-Net baseline. These efficiency gains enable deployment on consumer-grade hardware (NVIDIA RTX 2060, 6GB VRAM), representing a concrete step towards democratising AI-assisted diagnostics in resource-constrained clinical settings.

The principal contributions of this work are fourfold. First, a superpixel-based graph construction pipeline that exploits inherent tumour sparsity to achieve 890$\times$ dimensionality reduction while preserving tumour-discriminative features across four MRI modalities. Second, empirical validation through systematic ablation studies that a 5-layer GraphSAGE architecture is optimal — additional depth (6 layers) increases parameter count by 30\% without improving accuracy. Third, the empirical discovery of batch size sensitivity in medical imaging GNN tasks, where batch size 32 significantly outperforms larger configurations (48, 64), providing actionable guidance for future work in this domain. Fourth, a rigorous validation framework comprising patient-level stratified 5-fold cross-validation with 15 automated integrity checks, establishing zero data leakage and full reproducibility.

Returning to the central research question — \textit{Can a graph-based neural network achieve competitive segmentation accuracy ($>$90\% Dice) while providing efficiency gains sufficient for resource-constrained clinical deployment?} — this thesis provides an affirmative answer. The 92.92\% ensemble Dice, combined with 12.7ms inference time, 2.1GB memory footprint, and 439K parameter count, demonstrates that competitive accuracy and practical efficiency are not mutually exclusive. As healthcare AI scales to serve diverse global populations, efficiency-aware architectures of this kind will be critical for ensuring that advanced diagnostic tools remain accessible regardless of hardware infrastructure.


\subsection{Limitations}

Several limitations of this work are acknowledged. First, graph construction (SLIC superpixel segmentation and feature extraction) adds approximately 2--3 seconds per patient as a preprocessing step, which limits true end-to-end real-time performance when graphs are not pre-computed. Second, the 15 handcrafted node features require domain expertise in MRI physics and tumour biology; they may not generalise without modification to tumour types outside the glioma spectrum represented in BraTS. Third, this work addresses binary tumour segmentation only, foregoing the clinically relevant sub-region delineation (Enhancing Tumour, Tumour Core, Whole Tumour) evaluated in standard BraTS benchmarks. Fourth, validation is limited to a single dataset (BraTS 2021); performance on different scanner protocols, patient populations, or tumour subtypes (e.g.\ meningioma, metastases) is untested. Fifth, a 1--2\% Dice gap remains below the best transformer-based methods on the Whole Tumour metric. Sixth, clinical integration barriers — including regulatory approval (FDA 510(k) / CE marking), hospital IT infrastructure requirements, and prospective radiologist validation studies — remain unaddressed and constitute necessary steps before clinical translation.


\subsection{Future Work and Directions}

Several directions emerge naturally from this work. Architecturally, incorporating attention mechanisms (e.g.\ graph attention layers) or hierarchical multi-scale graph representations could close the remaining 1--2\% gap with transformer-based methods without sacrificing efficiency. A multi-task extension enabling simultaneous binary and multi-class segmentation (ET, TC, WT) would bring the approach to full BraTS benchmark comparability and significantly increase clinical utility.

On the preprocessing side, GPU-accelerated superpixel segmentation or differentiable clustering approaches could reduce the 2--3 second graph construction overhead, enabling true end-to-end real-time processing. Federated learning extensions would leverage the lightweight model architecture to enable privacy-preserving multi-institutional training without sharing patient data.

For clinical translation, the immediate priorities are prospective radiologist evaluation studies, multi-institutional validation on diverse scanner hardware, and regulatory submission preparation. Application of the same graph-based pipeline to other anatomical regions (prostate, liver, lung) and to longitudinal tumour tracking for treatment response assessment are further avenues of investigation.

This thesis establishes graph neural networks as a compelling paradigm for efficient brain tumour segmentation, demonstrating that competitive medical image analysis does not require massive volumetric models. As healthcare AI scales to serve global populations, efficiency-aware architectures will be critical for ensuring that advanced diagnostic tools are accessible across all levels of infrastructure and resource availability.